\documentclass[a4paper,11pt]{article}
\usepackage[utf8]{inputenc}
\usepackage[T1]{fontenc}
\usepackage[french]{babel}
\usepackage[left=2.5cm,right=2.5cm,top=2cm,bottom=2cm]{geometry}
\usepackage{hyperref}
\usepackage{xcolor}
\usepackage{fontawesome5}
\usepackage{lmodern}
\usepackage{setspace}
\usepackage{graphicx}

% --- Couleurs ---
\definecolor{primary}{RGB}{35, 55, 123}

% --- Configuration des liens ---
\hypersetup{
    colorlinks=true,
    linkcolor=primary,
    urlcolor=primary,
}

\begin{document}

% --- EN-TÊTE ---
\begin{flushleft}
\textbf{Loïc Aron MBASSI EWOLO} \\
Élève-Ingénieur 2A - ENSEA \\
\faPhone\ (+33) 7 59 52 33 98 \\
\faEnvelope\ loicaron.mbassiewolo@ensea.fr \\
\faLinkedin\ \href{https://www.linkedin.com/in/loic-aron-mbassi-ewolo-3942a9246/}{loic-mbassi} \\
\end{flushleft}

\vspace{0.5cm}

\begin{flushright}
Cergy, le 26 janvier 2026
\end{flushright}

\vspace{0.5cm}

\begin{flushleft}
\textbf{À l'attention de la Commission des Bourses} \\
ENSEA - École Nationale Supérieure de l'Électronique \\
et de ses Applications \\
6 Avenue du Ponceau \\
95014 Cergy-Pontoise
\end{flushleft}

\vspace{0.8cm}

\begin{center}
\textbf{\large Lettre de Motivation - Candidature à la Bourse d'Excellence aux Internationaux}
\end{center}

\vspace{0.5cm}

\onehalfspacing

Madame, Monsieur,

\vspace{0.3cm}

Je me permets de vous soumettre ma candidature pour la bourse d'excellence destinée aux étudiants internationaux en double diplôme de l'ENSEA. Actuellement en deuxième année du cycle ingénieur dans le cadre d'un double diplôme avec l'École Polytechnique de Yaoundé (Cameroun), je souhaite par cette lettre vous exposer ma situation, mon parcours académique, mon projet professionnel ainsi que mon implication dans la vie de l'école.

\vspace{0.4cm}

\textbf{\color{primary}Mon parcours académique et mes résultats}

\vspace{0.2cm}

À l'École Polytechnique de Yaoundé, j'ai obtenu une moyenne générale pondérée de \textbf{2,91/4 (soit 14,03/20)}, me classant \textbf{20\textsuperscript{ème} sur 108 étudiants} de ma promotion en Génie Informatique. Ce parcours m'a permis d'acquérir de solides bases en mathématiques appliquées, génie logiciel et intelligence artificielle.

Depuis mon arrivée à l'ENSEA en 2025, je poursuis avec sérieux et implication ma formation depuis le premier semestre et je suis maintenant en option Drones. Malgré les difficultés financières que je traverse, je maintiens un niveau académique satisfaisant : \textbf{je n'ai aucun rattrapage pour ce semestre 7}. Ce résultat témoigne de ma détermination à réussir, même dans des conditions difficiles.

\vspace{0.4cm}

\textbf{\color{primary}Ma situation financière actuelle}

\vspace{0.2cm}

Je tiens à être transparent sur ma situation financière, qui constitue un obstacle majeur à mon épanouissement académique. N'ayant pas de soutien financier familial stable (mes parents ayant organisé mon accueil via une famille d'accueil en France sans pouvoir assurer le financement de mes études), je dois subvenir seul à mes besoins.

Pour financer ma scolarité et mon quotidien, \textbf{je donne des cours particuliers chez Complétude}. Bien que cette activité me permette de générer des revenus (dont je peux fournir les relevés bancaires), elle représente une charge de travail conséquente qui \textbf{empiète sur mon temps d'étude et m'épuise physiquement et mentalement}. De plus, j'ai contracté une \textbf{dette de 2~000~€ auprès des alumni de l'ENSEA} pour pouvoir couvrir des frais urgents et \textbf{le reste de ma scolarité (4~231~€)}, dette que je dois également rembourser.

Cette situation m'oblige à jongler constamment entre mes obligations académiques et mes contraintes financières, ce qui limite ma capacité à donner le meilleur de moi-même dans mes études.

\vspace{0.4cm}

\textbf{\color{primary}Mon projet professionnel}

\vspace{0.2cm}

Mon objectif est de devenir \textbf{Ingénieur R\&D spécialisé en Intelligence Artificielle Embarquée}. Je suis particulièrement passionné par l'intersection entre les systèmes embarqués et le deep learning, domaine qui présente des défis fascinants en termes d'optimisation et de déploiement de modèles sur des plateformes à ressources limitées.

Cette passion s'est concrétisée par plusieurs expériences significatives :
\begin{itemize}
    \item \textbf{Stage de recherche au MILA} (Institut Québécois d'IA, été 2025) : travaux sur le phénomène de « Grokking » et l'interprétabilité des réseaux de neurones.
    \item \textbf{Projet UROP 2024} : développement d'un modèle de détection d'anomalies cardiaques à partir d'images d'ECG, en collaboration avec des mentors de Google DeepMind.
    \item \textbf{Challenge CoHoMa} (en cours ici à l'ENSEA) : optimisation d'algorithmes de SLAM et déploiement d'IA embarquée (YOLOv10) sur Nvidia Jetson pour la robotique autonome militaire.
\end{itemize}

À court terme, je recherche un stage de fin d'études (avril-septembre 2026) dans le domaine de l'IA embarquée ou du MLOps, puis un contrat de professionnalisation pour ma dernière année (j'ai déjà eu deux entretiens qui ne se sont malheureusement pas conclus par un contrat).

\vspace{0.4cm}

\textbf{\color{primary}Mon implication dans la vie de l'école}

\vspace{0.2cm}

\underline{À l'ENSEA :}
\begin{itemize}
    \item Je suis \textbf{membre actif de l'association ARES}.
    \item Je souhaite \textbf{participer à la Coupe de France de Robotique}, événement emblématique qui me permettrait d'appliquer mes compétences en systèmes embarqués et en IA.
\end{itemize}

\underline{À l'École Polytechnique de Yaoundé}, j'ai démontré un engagement associatif fort :
\begin{itemize}
    \item \textbf{Président du Club Informatique} et \textbf{Chef de la Cellule Projet}.
    \item \textbf{Lead de la Cellule IA} : j'ai contribué à augmenter le nombre de membres de plus de 50\% et organisé des collaborations avec Google DeepMind.
    \item \textbf{Animation de workshops techniques} : formations en Laravel, React, Machine Learning et Fine-tuning de LLM pour mes camarades.
    \item \textbf{Participation à de nombreux hackathons} avec des résultats notables :
    \begin{itemize}
        \item 1\textsuperscript{ère} place - IndabaX Africa 2025 (IA pour la Santé)
        \item 1\textsuperscript{ère} place - CITS National Hackathon (Urban Waste Management)
        \item 1\textsuperscript{ère} place - Mountain Hub Regional 2024
        \item 2\textsuperscript{ème} place - Semaine de la Jeune Fille Scientifique (Violence Detection)
    \end{itemize}
    \item \textbf{Volontaire à IndabaX Africa} et rédacteur d'articles techniques sur Medium.
\end{itemize}

\vspace{0.4cm}

\textbf{\color{primary}Pourquoi cette bourse est essentielle pour moi}

\vspace{0.2cm}

L'obtention de cette bourse d'excellence me permettrait de :
\begin{enumerate}
    \item \textbf{Réduire significativement ma charge de travail extrascolaire}, me libérant du temps et de l'énergie pour me consacrer pleinement à mes études et améliorer mes résultats académiques.
    \item \textbf{Rembourser progressivement ma dette} envers les alumni, allégeant ainsi ma charge mentale.
    \item \textbf{M'investir davantage dans la vie associative de l'ENSEA}, notamment dans le projet de la Coupe de Robotique avec ARES.
    \item \textbf{Donner le meilleur de moi-même} pour représenter dignement l'ENSEA et honorer la confiance que l'école m'accorde.
\end{enumerate}

\vspace{0.4cm}

Je suis convaincu que mon parcours, ma détermination et mon engagement démontrent ma capacité à surmonter les obstacles et à exceller académiquement lorsque les conditions le permettent. Cette bourse serait un \textbf{investissement dans un étudiant motivé}, passionné et déterminé à réussir.

\vspace{0.3cm}

Je reste à votre entière disposition pour tout entretien ou complément d'information.

\vspace{0.3cm}

Je vous prie d'agréer, Madame, Monsieur, l'expression de mes salutations respectueuses.

\vspace{0.5cm}

\begin{flushright}
\includegraphics[height=2cm]{signature.png} \\
Loïc Aron MBASSI EWOLO
\end{flushright}

\vspace{0.5cm}

\noindent\rule{\textwidth}{0.4pt}

\vspace{0.3cm}

\textbf{Pièces jointes :}
\begin{itemize}
    \item Relevé de notes de l'École Polytechnique de Yaoundé (avec classement)
    \item Relevés bancaires attestant de mes revenus (cours Complétude)
    \item Attestation de dette auprès des alumni ENSEA
\end{itemize}

\end{document}